\documentclass[10pt, a4paper]{article}
\usepackage[T2A]{fontenc}
\usepackage[utf8]{inputenc}
\usepackage{enumitem}
\usepackage[russian]{babel}
\usepackage{tikz}
\usepackage{amsmath, amsfonts, amssymb, amsthm, tikz}
\usepackage{cancel}
\usepackage{geometry}
\geometry{top=1.5cm, bottom=1.5cm, left=1.5cm, right=1.5cm}
\usetikzlibrary{decorations.pathreplacing, arrows.meta}



\title{\vspace{-1cm}\textbf{Конспект по дискретной математике}\vspace{-0.5cm}}
\author{}
\date{}

\def\divby{%
  \mathrel{\text{\vbox{\baselineskip.65ex\lineskiplimit0pt\hbox{.}\hbox{.}\hbox{.}}}}%
}

\begin{document}
\maketitle
\vspace{-0.8cm}

\section*{Список используемых обозначений}
\begin{tabular}{ll}
    $\lfloor x \rfloor$ & --- Округление числа $x$ вниз \\
    $\lceil x \rceil$ & --- Округление числа $x$ вверх \\
    $\iff$ & --- Тогда и только тогда, когда \\

\end{tabular}

\vspace{0.5cm}
\hrule

\section*{\centerline{\Large\textbf{\S \,\, Методы доказательств}}}
\vspace{0.4cm}

\subsection*{Конструктивное док-во (пример)}

\textbf{Пример:} Докажите, что пифагоровых троек бесконечно много.
\[ 3^2 + 4^2 = 5^2 \]
\[ (3k)^2 + (4k)^2 = (5k)^2 \text{ --- бесконечная серия, } \forall k \in \mathbb{N} \]

\subsection*{Контрпример}

\textbf{Вопрос:} Верно ли, что все слова русского языка не содержат более 4 согласных подряд? Нет, «контрпример» имеет 5.

\subsection*{Прямое док-во}

Если $a \divby c$ и $b \divby d$, то $ab \divby cd$.
\[ \left.
\begin{aligned}
a \divby c &\implies \exists k: a = ck \\
b \divby d &\implies \exists l: b = dl
\end{aligned}
\right\} \implies ab = ckdl \implies ab \divby cd \text{ по опр.} \]

\subsection*{Док-во от противного}
\noindent \hfill \[ A \longrightarrow B \iff \overline{B} \longrightarrow \overline{A} \]
\noindent
\textbf{Пример:} Доказать, что простых чисел бесконечно много.
\begin{proof}[Док-во]
\noindent
Пусть не так, тогда $p_1, \dots, p_n$ --- все простые числа. \\
Рассмотрим $x = p_1 \cdot p_2 \cdot p_3 \cdot \dots \cdot p_n + 1$, заметим, что 
$\left.
\begin{aligned}
x > p_i, \; \forall i \\
x \divby p_i, \; \forall i
\end{aligned}
\right\} \implies x \text{ --- новое простое число.}$
\end{proof}

\subsection*{Оценка + пример}

\textbf{Вопрос:} Какое наибольшее кол-во не бьющих друг друга ладей?

\begin{proof}[Док-во]
Пример на 8:
\begin{center}
\begin{picture}(80,80)
\multiput(0,0)(0,10){9}{\line(1,0){80}}
\multiput(0,0)(10,0){9}{\line(0,1){80}}
\multiput(5,5)(10,10){8}{\circle*{4}}
\end{picture}
\end{center}

\noindent От противного: Пусть 9, тогда 2 будут на одной горизонтали и будут бить друг друга. \\
\textit{На самом деле здесь мы пользуемся принципом Дирихле.}
\end{proof}

\subsection*{Принцип Дирихле}

\textbf{Утв.} Даны $n$ клеток и $n+1$ кроликов, то существует клетка с $\ge 2$ кроликами.

\paragraph{Утв. Обобщенный принцип Дирихле} 
Даны $n$ клеток, $k$ кроликов, то:
\begin{enumerate}
    \item $\exists$ клетка с $\ge \lceil \frac{k}{n} \rceil$ кроликами.
    \item $\exists$ клетка с $\le \lfloor \frac{k}{n} \rfloor$ кроликами.
\end{enumerate}

\begin{proof}[Док-во обобщенного принципа Дирихле:]
\noindent 
\begin{enumerate}
    \item От противного: пусть в каждой $< \lceil \frac{k}{n} \rceil$. В клетках $a_1, \dots, a_n$ кроликов. \\
    \[ a_i < \left\lceil \frac{k}{n} \right\rceil \implies a_i \le \left\lceil \frac{k}{n} \right\rceil - 1 < \frac{k}{n} \]
    \[ k = \sum a_i < n \cdot \frac{k}{n} = k \quad ?! \]

    \item Пусть в каждой $> \lfloor \frac{k}{n} \rfloor$.
    \[ a_i > \left \lfloor \frac{k}{n} \right \rfloor \implies a_i \ge \left \lfloor \frac{k}{n} \right \rfloor + 1 > \frac{k}{n} \]
    \[ k = \sum a_i > n \cdot \frac{k}{n} = k \quad ?! \]
\end{enumerate}
\end{proof}

\textbf{Пример:} 40 человек. Докажите, что есть месяц с $\ge 4$ днями рождения и есть месяц с $\le 3$. \\

\begin{proof}{\textit{по принципу Дирихле:}}
\begin{enumerate}
    \item Пусть в каждом $\le 3 \implies$ всего дней рожд $\le 36$, но людей $40$?!
    
    \item Пусть в каждом $\ge 4 \implies$ всего $\ge 48$, но людей $40$?!
\end{enumerate}
\end{proof}

\paragraph{Утв. (Бесконечный принцип Дирихле)} 
Есть бесконечное число кроликов и конечное число клеток, тогда $\exists$ клетка, в которой бесконечное число кроликов.
\begin{proof}[Док-во:]
Пусть нет, тогда в каждой клетке конечное число кроликов, но $\text{конечное} \times \text{конечное} = \text{конечное} < \text{бесконечное}$.
\end{proof}

\subsection*{Комбинаторные док-ва}

Вместо счета объектов и проверки их равенства, можно установить соответствие.

\subsection*{Подсчет двумя способами}
Часто используется в док-ве от противного: «что-то одновременно больше и меньше нуля».

\subsection*{Принцип крайнего}

Хорошая штука, когда не хватает условий. Часто работает с методом от противного, но это не обязательно.
\noindent        
\textbf{Задача:} Докажите, что среди $n\ge 3$ точек на плоскости есть такая, что все углы с другими точками $<180^\circ$. 
\begin{proof}[Док-во:]
    Рассмотрим самую левую точку и из таких самую нижнюю, заметим, что она подойдет, так как нет точек левее ее и $180^\circ$ тоже не бывает. 
\end{proof}

\subsection*{Инвариант --- «что-то, что не меняется»}

\textbf{Задача:} На доске написано число $a$. Разрешается его стереть и написать либо $3a$, либо $a+8$, либо $\frac{a}{5}$, если $a \divby 5$. Можно ли из $1$ получить $2$?

\noindent \textit{Док-во:} Нет, т.к. число всегда нечётное. Инвариант --- чётность.

\noindent \textbf{Вопрос:} А из числа $2$ в $4$? Нет, но инвариант сложнее.

\subsection*{Вспомогательная раскраска}

\textbf{Задача:} Из шахматной доски вырезали 2 противоположных угла. Можно ли её разрезать на доминошки $1 \times 2$?

\noindent \textit{Док-во:} Отметим шахматную раскраску. Пусть отрезанные углы чёрные, тогда у нас 32 белых, 30 чёрных, а каждая доминошка содержит по одной клетке каждого цвета $\implies$ не выйдет.

\noindent Инвариант: кол-во белых $-\,\, \text{кол-во}$ чёрных.

\subsection*{Полуинвариант --- «что-то, что меняется лишь в одну сторону»}

\textbf{Пример:} предыдущая задача, но можно ещё вырезать чёрные клетки как действие.

\noindent Полуинвариант: кол-во белых $-\,\, \text{кол-во}$ чёрных. Он только увеличивается, изначально 2, а должен стать 0 ?!

\subsection*{Индукция}

Нужна, чтобы доказывать утверждение для всех $n$ ($\forall n$).\\

\noindent \textbf{Пример 1:} $1 + 2 + 3 + \dots + n = \frac{n(n+1)}{2}$

\noindent База: $n = 1 \implies 1 = \frac{1(1+1)}{2}$

\noindent Переход: $n = k \implies n = k+1$. Хотим для $k+1$:
\[ 1 + 2 + 3 + \dots + k + (k+1) = \underbrace{\frac{k(k+1)}{2}}_{\text{= по ИП}} + (k+1) = \frac{(k+1)(k+2)}{2} = \frac{(k+1)((k+1)+1)}{2} \quad \checkmark \]

\subsection*{Возвратная индукция --- переход из $1, \dots, n \implies n+1$}

\textbf{Пример:} Алгоритм Евклида заканчивается. Делаем индукцию по $k = n + m$, переход от всех меньших k к k.

\noindent Переход: $(n, m) = (n - m, m)$ --- для второй пары заканчивается по ИП (индукционному предположению), т.к. $(n - m) + m < n + m$.

\subsection*{Переход: $n \implies n+x$}

\textbf{Задача:} Какую сумму можно выдать монетами номинала 3 и 5? \\
Докажем, что любую сумму $n \ge 8$ можно выдать.

\noindent Переход: $n \implies n+3$. Выдадим 3, а $n$ --- по ИП.

\noindent База: 
$n = 8 = 3 + 5; \quad n = 9 = 3 + 3 + 3; \quad n = 10 = 5 + 5$.

\subsection*{Переход: $2^n \implies 2^{n+1}$ и $x \implies x-1$}

\textbf{Задача:} Доказать, что любое число раскладывать в сумму степеней двойки, где каждая не более одного раза.

\noindent Переход: $n \implies 2n$ и $n \implies 2n+1$.

{\color{gray}V: $\checkmark$}

\section*{\centerline{\Large\textbf{\S \,\, Введение в теорию множеств}}}
\vspace{0.4cm}

Множество — первичное математическое понятие, которому не дано строгое математическое определение. Множество состоит из элементов. Множество задается набором своих элементов.

\vspace{0.2cm}
\noindent \textbf{Обозн.} $A$ --- множество, $x$ --- элемент, $x \in A$.

\vspace{0.4cm}

\noindent \textbf{Опр.}
Множество $A$ называется подмножеством множества $B$, если $\forall x \in A \implies x \in B$.

\vspace{0.2cm}
\noindent \textbf{Обозн.} $A \subset B$.

\vspace{0.2cm}
\noindent \textbf{Замечание} Иногда используют символ $\subseteq$, что означает нестрогое включение, т.е. множество \(A\) является подмножеством \(B\) и может быть равным ему (\(A = B\)). Если $A \neq B$, то $A$ называется собственным подмножеством и обозначается $A \subsetneq B$.

\vspace{0.4cm}

\noindent \textbf{Утв.}
Если $A \subset B$ и $B \subset A$, то $A = B$.

\vspace{0.2cm}
\begin{proof}[Док-во:]
$\forall x \in A \implies x \in B \quad \text{и} \quad \forall x \in B \implies x \in A$
\end{proof}

\noindent \textbf{Опр.}
Множество, не содержащее ни одного элемента, называется пустым.

\vspace{0.2cm}
\noindent \textbf{Обозн.} $\varnothing$

\vspace{0.6cm}

\noindent \textbf{Замечание (Парадокс Рассела / Брадобрея)}

\vspace{0.2cm}
\noindent Рассмотрим множество всех множеств, не содержащих самих себя. 
Вопрос: содержит ли это множество само себя?

\[ A = \{ B \mid B \notin B \} \]
\newpage
\noindent Разберем два возможных случая:
\begin{itemize}
    \item Если $A \in A$, то по определению множества $A \implies A \notin A$.
    \item Если $A \notin A$, то по определению множества $A \implies A \in A$ \quad ?!
\end{itemize}

На самом деле, нельзя так задавать множества. Нужно писать:
\[ A = \{ B \in \Omega \mid \text{условие на } B \} \]

\noindent Поэтому везде будем считать, что все множества лежат в некотором большом множестве $\Omega$.

\vspace{0.3cm}

\subsection*{Как задавать множества?}

\begin{enumerate}
    \item Поэлементно, например $A = \{3, \text{стул}, \text{мем с котиком}\}$
    \item Через условия, которые называются предикатами. \\
    $A = \{ x \in \mathbb{N} \mid x \;\divby\; 5 \}$ \\
    $P(x) = x \;\divby\; 5$ --- предикат.
    
    \textit{Замечание.} Не любое условие является предикатом, далее мы будем это игнорировать
\end{enumerate}

\vspace{0.3cm}
\noindent \textbf{Пример (Парадокс Берри)}

\vspace{0.2cm}
\noindent Рассмотрим множество $A = \{ x \in \mathbb{N} \mid x \text{ можно описать менее, чем 12 словами русского языка} \}$.

\vspace{0.2cm}
\noindent С одной стороны, множество $A$ конечно. Тогда существует наименьшее число $y$, которое в нем не лежит. Но тогда: $y$ --- это «наименьшее число, которое нельзя описать менее, чем двенадцатью словами русского языка» $\implies y \in A$ \quad ?!

\vspace{0.3cm}
\subsection*{Операции над множествами}

\begin{itemize}
    \item Множество $A \cap B = \{ x \mid x \in A \land x \in B \}$ называется \textbf{пересечением} ($\land$ --- конъюнкция --- логическая связка «и»).
    \item Множество $A \cup B = \{ x \mid x \in A \lor x \in B \}$ называется \textbf{объединением} ($\lor$ --- дизъюнкция --- логическая связка «или»).
\end{itemize}

\vspace{0.2cm}
\noindent \textbf{Замечание} Формально нужно писать $\{ x \in \Omega \mid \dots \}$, но в наивной теории множеств это подразумевается.

\begin{itemize}
    \item Множество $\overline{A} = \{ x\mid  x \in \Omega \land x \notin A \}$ называется \textbf{дополнением}
    \item Множество $A \setminus B = \{ x \mid x \in A \land x \notin B \}$ называется \textbf{разностью}
    \item Множество $A \Delta B = (A \setminus B) \cup (B \setminus A)$ называется \textbf{симметрической разностью} \\
    \small (множество элементов, лежащих ровно в одном из множеств)
\end{itemize}

\vspace{0.2cm}
\subsection*{Свойства операций над множествами:}

\begin{itemize}
    \item \textbf{Коммутативность:} 
    $A \cap B = B \cap A, \quad A \cup B = B \cup A, \quad A \Delta B = B \Delta A$
    
    \vspace{0.2cm}
    \begin{proof}[Док-во на примере объединения:]
    \[ A \cup B = \{ x \in \Omega \mid x \in A \lor x \in B \} \]
    \[ B \cup A = \{ x \in \Omega \mid x \in B \lor x \in A \} \]
    $x \in A \lor x \in B \iff x \in B \lor x \in A$ (т.к. логическая дизъюнкция коммутативна)
    \end{proof}

    \vspace{0.1cm}
    \item \textbf{Ассоциативность:}
    $A \cap (B \cap C) = (A \cap B) \cap C$ \\
    $A \cup (B \cup C) = (A \cup B) \cup C$ \\
    $A \Delta (B \Delta C) = (A \Delta B) \Delta C$
    \begin{proof}[Док-во:]
    \[ A \Delta B = \{ x \mid x \in A \text{ xor } x \in B \} \]
    Логическое исключающее «или» (xor) ассоциативно, откуда следует ассоциативность для любого числа множеств.
    \end{proof}

    \vspace{0.3cm}
    \item \textbf{Дистрибутивность:} \\
    $A \cup (B \cap C) = (A \cup B) \cap (A \cup C)$ \\
    $A \cap (B \cup C) = (A \cap B) \cup (A \cap C)$
    
    \vspace{0.2cm}
    \begin{proof}[Док-во:]
    \[ A \cup (B \cap C) = \{ x \mid x \in A \lor (x \in B \land x \in C) \} \]
    \[ x \in A \lor (x \in B \land x \in C) \iff (x \in A \lor x \in B) \land (x \in A \lor x \in C) \]
    Полученное выражение в точности совпадает с определением правой части:
    \[ (A \cup B) \cap (A \cup C) = \{ x \mid (x \in A \lor x \in B) \land (x \in A \lor x \in C) \} \]
    \end{proof}
    \item \textbf{Закон двойного отрицания:}
    $\overline{\overline{A}} = A$

    \item \textbf{Законы де Моргана:} $\overline{A \cap B} = \overline{A} \cup \overline{B}$, \quad $\overline{A \cup B} = \overline{A} \cap \overline{B}$
    \end{itemize}
    
\noindent Множество упорядоченных пар $A \times B = \{ (x, y) \mid x \in A \text{ и } y \in B \}$ называется \textbf{декартовым произведением}

\vspace{0.4cm}

\noindent \textbf{Опр.}
\textbf{(Бинарным) отношением} $R$ называется любое подмножество декартова произведения $A \times B$.

\vspace{0.2cm}
\noindent \textbf{Пример:}
Отношение нестрогого порядка $\le$ для натуральных чисел: \\
${\le} = \{ (1, 1), (1, 2), (1, 3), (2, 2), (2, 3), \dots \}$ 

\vspace{0.2cm}
\noindent Вместо записи $(x, y) \in R$ пишут: $xRy$.

\vspace{0.4cm}
\noindent \textbf{Опр.}
\textbf{Функция} --- это такое отношение $R$, которое обладает следующими свойствами:
\[ \begin{cases} 
    \forall x \in A \;\; \exists y \in B : (x, y) \in R \\ 
    \forall x \in A \;\; \forall y_1, y_2 \in B : \big( (x, y_1) \in R \land (x, y_2) \in R \big) \implies y_1 = y_2 
\end{cases} \]

\vspace{0.4cm}

\noindent \textbf{Опр.}
Два множества называются \textbf{равномощными}, если между ними существует взаимно однозначное соответствие (\textbf{биекция}).

\vspace{0.4cm}

\noindent \textbf{Опр.}
Множество $A$ \textbf{не больше по мощности}, чем множество $B$, если существует соответствие элементов $B$ элементам $A$ (но не обязательно всем элементам $B$), являющееся \textbf{инъекцией}.

\vspace{0.2cm}
\noindent \textbf{Обозн.} $|A| = |B| \iff A \sim B$ \quad $|A| \le |B| \iff A \preceq B$

\vspace{0.4cm}

\noindent \textbf{Пример:} Интуитивно кажется, что $|\mathbb{N}| \prec |\mathbb{Z}|$, но на самом деле их мощности равны: $|\mathbb{N}| = |\mathbb{Z}|$.

Сопоставим каждому целому числу натуральное по формуле 
\begin{equation*}
    f(x) = \begin{cases} 
2x+1, & \text{если } x \ge 0, \\ 
2|x|, & \text{если } x < 0. 
\end{cases}
\end{equation*}

Заметим, что оно взаимно однозначное.

\vspace{0.4cm}

\noindent \textbf{Опр.}
Множество называется \textbf{счетным}, если оно равномощно множеству натуральных чисел $\mathbb{N}$ (т.е. все его элементы можно пронумеровать).

\vspace{0.4cm}

\noindent $\mathbb{Q}$ --- счётное множество, $\frac{p}{q} \in \mathbb{Q}, \quad p \in \mathbb{Z}, \quad q \in \mathbb{N}$


\begin{center}
\begin{tikzpicture}[scale=1.2, >=stealth]
    \draw[->, thick] (-3.5,0) -- (3.5,0) node[right] {$p \in \mathbb{Z}$};
    \draw[->, thick] (0,-0.5) -- (0,4.5) node[above] {$q \in \mathbb{N}$};
    
    
    \foreach \p in {-3,-2,-1,0,1,2,3} {
        \foreach \q in {1,2,3} {
            \filldraw[black] (\p,\q) circle (1.5pt);
            \ifnum\p=0
                \node[below right, scale=0.65, text=gray] at (\p,\q) {$0$};
            \else
                \node[below right, scale=0.65, text=gray] at (\p,\q) {$\frac{\p}{\q}$};
            \fi
        }
    }
    
    \draw[->, red, thick] (0,1) -- (1,1);
    \draw[->, red, thick] (1,1) -- (0,2);
    \draw[->, red, thick] (0,2) -- (-1,1);
    \draw[->, red, thick] (-1,1) -- (-2,1);
    \draw[->, red, thick] (-2,1) -- (-1,2);
    \draw[->, red, thick] (-1,2) -- (0,3);
    \draw[->, red, thick] (0,3) -- (1,2);
    \draw[->, red, thick] (1,2) -- (2,1);
    \draw[->, red, thick] (2,1) -- (3,1);
    \draw[->, red, thick] (3,1) -- (2,2);
    \draw[->, red, thick] (2,2) -- (1,3);
    \draw[->, red, thick] (1,3) -- (0,4);
    \draw[->, red, thick] (0,4) -- (-1,3);
    \draw[->, red, thick] (-1,3) -- (-2,2);
    \draw[->, red, thick] (-2,2) -- (-3,1) node[left, scale=0.8, text=red] {$\dots$};

    % Метки на горизонтальной оси p
    \foreach \x in {-3,-2,-1,1,2,3} \node[below=1mm, scale=0.8] at (\x,0) {\x};
    \node[below left=1mm, scale=0.8] at (0,0) {0};
    
    % Метки на вертикальной оси q
    \foreach \y in {1,2,3,4} \node[left=1mm, scale=0.8] at (0,\y) {\y};
\end{tikzpicture}
\end{center}


\vspace{0.4cm}

\noindent \textbf{Теорема}
Отрезок $[0, 1]$ --- несчётное множество.

\vspace{0.3cm}

\begin{proof}[Доказательство (Диагональный метод Кантора):]
Пусть мы занумеровали все числа отрезка. Запишем их в десятичной записи:
\begin{align*}
    1. & \quad 0, \textbf{1}11111\dots \\
    2. & \quad 0,1\textbf{2}00\dots \\
    3. & \quad 0,10\textbf{1}001\dots \\
    4. & \quad 0,146\textbf{9}2\dots
\end{align*}

\noindent Построим число, которое отличается от первого в первом знаке после запятой, от второго --- во втором знаке после запятой и т.д. Будем брать любую цифру, кроме той, которая там стоит, и кроме цифры $9$, чтобы не возникло периодического переноса разряда в виде $0.999\dots = 1.0$.

\vspace{0.2cm}
\noindent Тогда это новое число отличается от всех чисел в списке. Следовательно, его в этом списке нет, что приводит к противоречию.
\end{proof}

\vspace{0.2cm}
\noindent \textbf{Опр.}
Мощность отрезка $[0, 1]$ называется \textbf{континуумом}.

{\color{gray}V: $\checkmark$}

\section*{\centerline{\Large\textbf{\S \,\, Начало комбинаторики}}}
\vspace{0.4cm}


\textbf{Опр.} Набор множеств $A_1, \dots, A_n$ называется \textit{покрытием} множества $X$, если:
$$A_1 \cup A_2 \cup \dots \cup A_n \supset X$$

\noindent \textbf{Опр.} Покрытие $A_1, \dots, A_n$ называется \textit{разбиением} множества $X$, если:
$$A_i \subset X \quad \text{и} \quad A_i \cap A_j = \varnothing \quad \forall i, j$$

\noindent \textbf{Обозн.} Если $A \cap B = \varnothing$, то вместо $A \cup B$ часто пишут $A \sqcup B$ (\textit{дизъюнктное объединение}).

Тогда разбиение множества $X$ можно записать в виде:
$$X = A_1 \sqcup A_2 \sqcup \dots \sqcup A_n$$

\subsection*{Правило суммы}

Для двух попарно непересекающихся множеств:
$$|A \sqcup B| = |A| + |B|$$

\textbf{Обобщенная версия (доказывается по мат. индукции):}
$$|A_1 \sqcup A_2 \sqcup \dots \sqcup A_n| = |A_1| + \dots + |A_n|$$


\vspace{0.3cm}
\noindent А если $A \cap B \neq \varnothing$, то как посчитать $|A \cup B|$?
$$|A \cup B| = |A| + |B| - |A \cap B|$$

\textbf{Док-во:}
$$A \cup B = A \cup (B \setminus A) = A \sqcup (B \setminus A) \implies |A \cup B| = |A| + |B \setminus A| = |A| + |B| - |A \cap B|$$

\noindent Часто, чтобы посчитать число элементов с каким-то условием, можно посчитать число всех остальных элементов. 

\textbf{Следствие.}
$$|\bar A| = |\Omega| - |A|$$



А если $A \cup B \cup C$? 

\vspace{0.2 cm}
Оказывается, формула такая:
$$|A \cup B \cup C| = |A| + |B| + |C| - |A \cap B| - |B \cap C| - |A \cap C| + |A \cap B \cap C|$$

\subsection*{Формула включений-исключений}

\[
\begin{aligned}
|A_1 \cup \dots \cup A_n|
={}& |A_1| + \dots + |A_n| - \underbrace{|A_1 \cap A_2| - \dots - |A_{n-1} \cap A_n|} _{\text{все пары мн-ва}} + \\
& + \underbrace{|A_1 \cap A_2 \cap A_3| + \dots + |A_{n-2} \cap A_{n-1} \cap A_n|} _{\text{все тройки мн-ва}} \dots \\
& \quad \dots \underset{\text{в зависимости от четности}}{\pm} |A_1 \cap A_2 \cap \dots \cap A_n|
\end{aligned}
\]

\textbf{Док-во:} По индукции

База: $n = 1, 2, 3$ $\checkmark$

Переход: $n \to n+1$

Рассмотрим мощность объединения $n+1$ множества:
$$|A_1 \cup A_2 \cup \dots \cup A_n \cup A_{n+1}|$$

Обозначим: $A_1 \cup \dots \cup A_n = X$

Тогда:
\[
\begin{aligned}
|X \cup A_{n+1}| ={} & |X| + |A_{n+1}| - |X \cap A_{n+1}| = \overbrace{|A_1| + \dots + |A_n| - |A_1 \cap A_2| - \dots - |A_{n-1} \cap A_n| + \dots}^{\text{все наборы из } n \text{ множеств без } A_{n+1}} \\
& \quad + |A_{n+1}| - |X \cap A_{n+1}|
\end{aligned}
\]


Как посчитать $|X \cap A_{n+1}|$?
$$X \cap A_{n+1} = (A_1 \cup \dots \cup A_n) \cap A_{n+1} =$$
$$= (A_1 \cap A_{n+1}) \cup (A_2 \cap A_{n+1}) \cup \dots \cup (A_n \cap A_{n+1})$$

Обозначим $A_i \cap A_{n+1}$ за $B_i$.

Применим ИП (индукционное предположение) для $B_i$, которых $n$ штук: 

\[
\begin{aligned}
|X \cap A_{n+1}| 
&= |B_1| + \dots + |B_n| - |B_1 \cap B_2| - \dots - |B_{n-1} \cap B_n| + \dots \\
&= |A_1 \cap A_{n+1}| + |A_2 \cap A_{n+1}| + \dots + |A_n \cap A_{n+1}| \\
&\quad- |A_1 \cap A_2 \cap A_{n+1}| - \dots - |A_{n-1} \cap A_n \cap A_{n+1}| + \dots \\
&\quad \dots \pm |A_1 \cap A_2 \cap \dots \cap A_n \cap A_{n+1}|
\end{aligned}
\]


То есть в точности все наборы с $A_{n+1}$, кроме самого $A_{n+1}$. Таким образом получаем формулу все множества для $n+1$ с нужными знаками, поскольку перед $|X \cap A_{n+1}|$ стоит минус.


\subsection*{Правило произведения}

$$|A \times B| = |A| \cdot |B|$$

\textbf{Опр.} Упорядоченным набором (множеством) называется множество, где каждый элемент имеет вид $(x, n)$, где все $n \in \mathbb{N}$ различны и принимают значения от $1$ до мощности множества.



\vspace{0.3cm}
\textbf{Замечание.}
Это формальная запись, упорядоченные наборы можно задать как угодно по-другому, но нужно как-то их задать, ведь просто так не выйдет.
$$\{x, y\} = \{y, x\}$$
$$\{\{x, 1\}, \{y, 2\}\} \neq \{\{y, 1\}, \{x, 2\}\}, \quad \text{если } x \neq y$$
\quad\quad Вместо $\{\{x, 1\}, \{y, 2\}\}$ пишут: $(x, y)$, и поэтому $(x, y) \neq (y, x)$

\vspace{0.3cm}
\textbf{Опр.} Мультимножеством $X$ называется множество, в котором «разрешены» повторения элементов.

Также как и с упорядоченным множеством их надо как-нибудь формально задать, например $(x,n)$, где $n \in \mathbb{N}$ будет означать $n$-ую копию элемента $x$.

\vspace{0.2cm}
\textbf{Замечание.}
Для обычных множеств верно, что
$$\{x, x\} = \{x\}$$

А мультимножества отличаются.
$$\{(x, 1), (x, 2)\} \neq \{(x, 1)\}$$

Тем не менее, почти всегда запись $\{a,a,b,b,b\}$ означает именно мультимножество с двумя элементами $a$ и тремя $b$. Для мультимножества нет стандартного обозначения.

\vspace{0.3cm}
\begin{tabular}{|l|c|c|}
\hline
 & \textbf{Упорядоченное (размещения)} & \textbf{Неупорядоченное (сочетания)} \\ \hline
\textbf{Без повторений} & 
\begin{tabular}[c]{@{}c@{}}(1,2), (1,3), (2,1),\\ (2,3), (3,1), (3,2)\end{tabular} & 
\{1,2\}, \{2,3\}, \{1,3\} \\ \hline
\textbf{С повторениями} & 
\begin{tabular}[c]{@{}c@{}}(1,1), (1,2), (1,3),\\ (2,1), (2,2), (2,3),\\ (3,1), (3,2), (3,3)\end{tabular} & 
\begin{tabular}[c]{@{}c@{}}\{1,2\}, \{2,3\}, \{1,3\},\\ \{1,1\}, \{2,2\}, \{3,3\}\end{tabular} \\ \hline
\end{tabular}

\paragraph{1) Размещения с повторениями}~\\[0.3em]
Пусть имеется $n$ различных объектов, и нужно выбрать $k$ из них (порядок важен, элементы могут повторяться). На каждое из $k$ мест можно выбрать любой из $n$ элементов. По принципу произведения --- это $n^k$.

\paragraph{2) Размещения без повторений}~\\[0.3em]
На первое место можно выбрать любой из $n$ элементов. На второе место --- любой, кроме одного уже выбранного --- $(n-1)$ способ. На третье место --- $(n-2)$ способов и т.д.

Значит, формула:
$$n(n-1)(n-2)\dots(n-k+1) = A_n^k$$

Частный случай при $k=n$ называется \textit{перестановкой} (размещения без повторений). 
Число перестановок (порядков):
$$P_n = n(n-1)\dots 2 \cdot 1 = n!$$

\textbf{Утв.} $$A_n^k = \frac{n!}{(n-k)!}$$

\textbf{Док-во (алгебраическое):}
$$\frac{n!}{(n-k)!} = \frac{n(n-1)(n-2)\dots(n-k+1)\cancel{(n-k)}\cancel{(n-k-1)}\dots \cancel{2\cdot 1}}{\cancel{(n-k)}\cancel{(n-k-1)}\dots \cancel{2\cdot 1}} = n(n-1)(n-2)\dots(n-k+1) = A_n^k$$

Можно доказать комбинаторно:

\vspace{0.2cm}
Среди всех порядков (перестановок) сколько тех, которые начинаются с фиксированного размещения из $k$?

Столько же, сколько способов упорядочить последние $n-k \implies$ Каждый способ посчитан $(n-k)!$ раз $\implies A_n^k = \frac{n!}{(n-k)!}$

\paragraph{3) Сочетания без повторений}~\\[0.3em]
Обозначим это число за $x$.
Тогда посчитаем число размещений без повторений двумя способами:
\begin{enumerate}
    \item С одной стороны, по определению: $A_n^k$
    \item С другой стороны, возьмем каждое сочетание без повторений и рассмотрим все его возможные порядки, которых $k!$, т.к. в каждом сочетании $k$ различных элементов
\end{enumerate}

Получаем: 
$$x \cdot k! = A_n^k$$
$$x = \frac{n(n-1)\dots(n-k+1)}{k!} = \frac{n!}{(n-k)!k!}$$

\textbf{Обозн.} $C_n^k$, иногда ${n \choose k}$

\textbf{Замечание} Так как $C_n^k \in \mathbb{N} \implies$ комбинаторно доказали, что $n(n-1)\dots(n-k+1) \ \divby \ k!$

$$C_n^k = \frac{n!}{k!(n-k)!}$$

\subsection*{Свойства $C_n^k$}

\noindent \textbf{1)} $C_n^k = C_n^{n-k}$

\vspace{0.2cm}
\noindent \textbf{Док-во:}
1) алгебраическое:
$$C_n^{n-k} = \frac{n!}{(n-k)!(n-(n-k))!} = C_n^k$$

\noindent 2) комбинаторное: Выбрать $k$ элементов — то же самое, что не выбрать $n-k$ элементов (выбрать те, которые не берем).

\vspace{0.2cm}

\noindent Можно это совсем строго записать через явное построение биекции: $|X| = n$, $A \subset X$, $|A| = k$.

\noindent $\text{Количество различных} A = C_n^k$

\noindent $A \longleftrightarrow X \setminus A$ — взаимно однозначное соответствие $\implies \text{количество различных}\ X \setminus A = C_n^k$

\vspace{0.2cm}
\noindent Но с другой стороны, $|X \setminus A| = n - k \implies \text{количество различных } X \setminus A = C_n^{n-k} \implies C_n^k = C_n^{n-k}$

\medskip
\noindent \textbf{2) Рекуррентное свойство}
$$C_{n+1}^{k+1} = C_n^{k+1} + C_n^k$$

\noindent \textbf{Док-во:}
1) алгебраическое:
$$C_{n+1}^{k+1} = \frac{(n+1)!}{(k+1)!(n-k)!}$$

$$C_n^k + C_n^{k+1} = \frac{n!}{k!(n-k)!} + \frac{n!}{(k+1)!(n-k-1)!} = \frac{n!}{k!(n-k-1)!} \left( \frac{1}{n-k} + \frac{1}{k+1} \right) =$$
$$= \frac{n!}{k!(n-k-1)!} \cdot \frac{n+1}{(n-k)(k+1)} = \frac{(n+1)!}{(k+1)!(n-k)!} = C_{n+1}^{k+1}$$

\vspace{0.2cm}
2) комбинаторное: Посчитаем число способов выбрать $k+1$ элемент из $n+1$ следующим образом. Зафиксируем какой-нибудь элемент.
\begin{itemize}
    \item Если мы его не берем, то нужно выбрать те же $k+1$ элементов, но теперь из $n$  элементов. Таких способов: $C_n^{k+1}$.
    \item Если же мы его берем, то теперь из $n$ элементов нужно выбрать оставшиеся $k$ штук (так как один уже выбран). Таких способов: $C_n^k$.
\end{itemize}

\noindent Заметим, что каждый из $C_{n+1}^{k+1}$ способов соответствует ровно одному из этих двух вариантов. Поэтому есть взаимно однозначное соответствие $\implies C_{n+1}^{k+1} = C_n^{k+1} + C_n^k$.

\vspace{0.2cm}
\vspace{0.2cm}
\noindent \textbf{Пример:} $X = \{1, 2, 3, 4\}$, ищем пары $C_4^2 = 6$: $\{1, 2\}, \{1, 3\}, \{1, 4\}, \{2, 3\}, \{2, 4\}, \{3, 4\}$. Зафиксируем «$4$» и разобьем пары на два множества:

\begin{minipage}[t]{0.55\textwidth}
    \centerline{\small множества с $4$}
    $$C_3^1 = 3$$
    $$\{1, 4\}, \{2, 4\}, \{3, 4\}$$
\end{minipage}
\hspace{0.5cm}
\begin{minipage}[t]{0.15\textwidth}
    \centerline{\small множества без $4$}
    $$C_3^2 = 3$$
    $$\{1, 2\}, \{1, 3\}, \{2, 3\}$$
\end{minipage}

\centerline{$C_4^2 = C_3^1 + C_3^2$}

\noindent \textbf{3) Тождество Вандермонда}
$$\sum_{k=0}^{r} C_m^k \cdot C_n^{r-k} = C_{m+n}^r$$

\noindent \textbf{Док-во (комбинаторное):} 
Нужно выбрать $r$ людей из группы, в которой $m$ мужчин и $n$ женщин.

\noindent С одной стороны, число способов — это $C_{m+n}^r$.

\noindent С другой стороны, посчитаем число способов, если в группе ровно $k$ мужчин. 
$$\text{Таких способов: } \underbrace{C_m^k}_{\text{выбрать } k \text{ мужчин}} \cdot \underbrace{C_n^{r-k}}_{\text{выбрать } r-k \text{ женщин}}$$

\noindent Заметим, что сумма по всем $k$ в точности дает все способы по одному разу:
$$\implies \sum_{k=0}^{r} C_m^k \cdot C_n^{r-k} = C_{m+n}^r$$

\subsection*{Треугольник Паскаля}

\noindent Заметим, что $C_n^0 = C_n^n = 1 \quad \forall n$, а также что $C_{n+1}^{k+1} = C_n^k + C_n^{k+1}$.

\vspace{0.2cm}
\noindent Рассмотрим треугольник:
$$
\begin{gathered}
1 \\
1 \quad 1 \\
1 \quad 2 \quad 1 \\
1 \quad 3 \quad 3 \quad 1 \\
1 \quad 4 \quad 6 \quad 4 \quad 1 \\
1 \quad 5 \quad 10 \quad 10 \quad 5 \quad 1
\end{gathered}
$$

\vspace{0.2cm}
\noindent Или в виде сочетаний:

\vspace{0.3cm}
\begin{center}
\begin{tikzpicture}[x=1.3cm,y=1.1cm, >=stealth, red_arrow/.style={->, red, thick, shorten <=2pt, shorten >=2pt}]
    \node (C00) at (0,0) {$C_0^0$};
    \node (C10) at (-0.5,-1) {$C_1^0$};
    \node (C11) at (0.5,-1) {$C_1^1$};
    
    \node (C20) at (-1,-2) {$C_2^0$};
    \node (C21) at (0,-2) {$C_2^1$};
    \node (C22) at (1,-2) {$C_2^2$};
    
    \node (C30) at (-1.5,-3) {$C_3^0$};
    \node (C31) at (-0.5,-3) {$C_3^1$};
    \node (C32) at (0.5,-3) {$C_3^2$};
    \node (C33) at (1.5,-3) {$C_3^3$};
    
    \node (C40) at (-2,-4) {$C_4^0$};
    \node (C41) at (-1,-4) {$C_4^1$};
    \node (C42) at (0,-4) {$C_4^2$};
    \node (C43) at (1,-4) {$C_4^3$};
    \node (C44) at (2,-4) {$C_4^4$};

    \node[red, font=\small] at (0,-1.3) {$+$};
    \node[red, font=\small] at (-0.5,-2.3) {$+$};
    \node[red, font=\small] at (0.5,-2.3) {$+$};
    \node[red, font=\small] at (-1,-3.3) {$+$};
    \node[red, font=\small] at (0,-3.3) {$+$};
    \node[red, font=\small] at (1,-3.3) {$+$};


    \draw[red_arrow] (C10) -- (C21);
    \draw[red_arrow] (C11) -- (C21);
    
    \draw[red_arrow] (C20) -- (C31);
    \draw[red_arrow] (C21) -- (C31);
    \draw[red_arrow] (C21) -- (C32);
    \draw[red_arrow] (C22) -- (C32);
    
    \draw[red_arrow] (C30) -- (C41);
    \draw[red_arrow] (C31) -- (C41);
    \draw[red_arrow] (C31) -- (C42);
    \draw[red_arrow] (C32) -- (C42);
    \draw[red_arrow] (C32) -- (C43);
    \draw[red_arrow] (C33) -- (C43);
\end{tikzpicture}
\end{center}
\vspace{0.3cm}

\noindent Из одинаковых начальных условий и рекуррентного свойства эти треугольники совпадают.

\vspace{0.2cm}
\noindent \textbf{Пример.} Есть прямоугольная доска $n \times m$ с фишкой в левом нижнем углу. Её можно двигать только вверх и вправо. Сколько различных (решетчатых) путей до правого верхнего угла?

\vspace{0.3cm}
\noindent
\begin{minipage}[c]{0.3\textwidth}
    \begin{tabular}{|c|c|c|c|c|}
    \hline
    1 & 3 & 6 & 10 & 15 \\ \hline
    1 & 2 & 3 & 4 & 5 \\ \hline
    1 & 1 & 1 & 1 & 1 \\ \hline
    \end{tabular}
\end{minipage}
\begin{minipage}[c]{0.6\textwidth}
    $\leftarrow$ количество способов дойти до данной клетки
\end{minipage}

\vspace{0.3cm}
\noindent Заметим, что количество способов дойти до данной клетки --- это просто сумма значений слева и снизу. Легко видеть, что это и есть треугольник Паскаля. Если доска $n \times m$, то ответ: $C_{(n-1)+(m-1)}^{m-1}$

\vspace{0.2cm}
\noindent \textbf{Докажем другим способом:}
Каждый путь длины ровно $(n-1) + (m-1)$, причем ровно $n-1$ шагов вправо и $m-1$ шагов вверх. Поэтом каждый путь соответствует расстановке $m-1$ шага вверх среди $(n-1) + (m-1)$ шагов пути, т.е. $C_{(n-1)+(m-1)}^{m-1}$.

\subsection*{Бином Ньютона}

$$(x+y)^n = \sum_{k=0}^{n} C_n^k \cdot x^{n-k} \cdot y^k = C_n^0 x^n + C_n^1 x^{n-1} y + C_n^2 x^{n-2} y^2 + \dots + C_n^{n-1} x^1 y^{n-1} + C_n^n y^n$$

\vspace{0.2cm}
\noindent \textbf{Док-во (комбинаторное):} 
$$(x+y)^n = \underbrace{(x+y)(x+y)(x+y)\dots(x+y)}_{n \text{ раз}}$$

\noindent Хотим посчитать, сколько раз встретится $x^{n-k} y^k$ при раскрытии скобок.
Встретится ровно столько раз, сколько способов выбрать $k$ штук $y$ из $n$ скобок, т.е. $C_n^k$

\vspace{0.3cm}
\noindent \textbf{Утв.} $C_n^0 + C_n^1 + \dots + C_n^n = 2^n$

\vspace{0.2cm}
\noindent \textbf{Док-во:} 
1) алгебраическое:
$(1+1)^n \stackrel{\text{по Б.Н.}}{=} C_n^0 \cdot 1^n \cdot 1^0 + C_n^1 \cdot 1^{n-1} \cdot 1^1 + C_n^2 \cdot 1^{n-2} \cdot 1^2 + \dots = \sum_{k=0}^{n} C_n^k \quad$

\medskip
\noindent 2) комбинаторное:

\noindent Посчитаем число результатов подбрасывания монетки $n$ раз. Их $2^n$. С другой стороны, тех, в которых ровно $k$ орлов — $C_n^k \implies \text{вся сумма } \sum_{k=0}^{n} C_n^k = 2^n \quad$

\vspace{0.3cm}
\noindent \textbf{Замечание:} Часто $C_n^k$ называют \textit{биномиальными коэффициентами}. Есть соглашение, что $C_n^k = 0$, если $k < 0$ или $k > n$.

\subsection*{Сочетания с повторениями}

\noindent Зафиксируем $X = \{1, \dots, n\}$

\noindent Хотим посчитать количество сочетаний с повторениями из $k$ элементов.

\vspace{0.2cm}
\noindent \textbf{Пример:} $X = \{1, 2, 3\}, \quad n = 3, \quad k = 2$

$$\{1, 1\}, \{1, 2\}, \{1, 3\}$$
$$\{2, 2\}, \{2, 3\}, \{3, 3\}$$

\noindent Зададим внутри каждого множества порядок по неубыванию и прибавим 1 ко второму числу:

$$\{1, 2\}, \{1, 3\}, \{1, 4\}$$
$$\{2, 3\}, \{2, 4\}, \{3, 4\}$$

\noindent Получили сочетания без повторов из элементов $\{1, \dots, n+k-1\}$, их число $C_{n+k-1}^k$ (для примера: $C_4^2$).

\vspace{0.2cm}
\noindent В общем случае в каждом множестве зафиксируем порядок также по неубыванию:
$$a_1 \le a_2 \le \dots \le a_k$$

\noindent Сопоставим ему строго возрастающий набор со сдвигом:
$$a_1 < a_2+1 < a_3+2 < \dots < a_k+k-1$$

\noindent Утверждается, что получились все сочетания без повторений $k$ элементов из $n+k-1$ элементов. Заметим, что все такие сочетания соответствуют какому-то сочетанию с повтором $\implies \text{их число } C_{n+k-1}^k$

\medskip
\noindent \textbf{Обозн.} $\left( \!\! \binom{n}{k} \!\! \right)$

\vspace{0.3cm}
\noindent \textbf{Сочетания с повторениями:} пусть каждый из $n$ элементов встречается $x_i$ раз, где $x_i \ge 0$.
Тогда мы ищем число целых неотрицательных решений уравнения:
$$ \sum_{i=1}^n x_i = k $$

\noindent \textbf{Переформулировка задачи («Шары и перегородки»):} есть $n$ различных ящиков (это наши элементы) и $k$ одинаковых шаров (это сколько раз мы их выбираем). Сколько способов рассовать шары по ящикам?

\vspace{0.2cm}
\noindent Представим это как строку из $k$ шаров и $n-1$ перегородок, разделяющих ящики. Всего позиций $n+k-1$. Выбираем $k$ мест для шаров:
$$C_{n+k-1}^k$$

% 6.10 (готово можно проверять)
\noindent
\textbf{Урновая схема}

\noindent Есть $n$ объектов в урне и вытаскиваются $k$ из них.

\begin{table}[h!]
\centering
\renewcommand{\arraystretch}{1.5}
\begin{tabular}{|l|c|c|}
\hline
\textbf{Предметы на входе} & \textbf{с возвращ.} & \textbf{без возвращ.} \\ \hline
\textbf{Упорядоч.}         & $n^k$               & $A_n^k$               \\ \hline
\textbf{Неупорядоч.}       & $\left(\!\binom{n}{k}\!\right)$ & $C_n^k$   \\ \hline
\end{tabular}
\end{table}

\noindent \textbf{Задача про разложение по ящикам}

\noindent Есть $n$ предметов, $k$ ящиков. Раскладываем предметы по ящикам.

\begin{table}[h!]
\centering
\renewcommand{\arraystretch}{1.5}
\begin{tabular}{|l|l|c|c|}
\hline
\textbf{Ящики} & \textbf{Предметы} & \textbf{\begin{tabular}[c]{@{}c@{}}Произвольное число\\ элементов в\\ каждом ящике\end{tabular}} & \textbf{\begin{tabular}[c]{@{}c@{}}Не более одного\\ предмета в ящике\end{tabular}} \\ \hline
различимые     & различимые        & $k^n$                                                                                        & $A_k^n$                                                                             \\ \hline
различимые     & неразличимые      & $\left(\!\binom{k}{n}\!\right)$                                                              & $C_k^n$                                                                             \\ \hline
\end{tabular}
\end{table}

\subsection*{Подсчёт отображений на конечных множествах}

\textbf{Обозн.} Множество функций из $X$ в $Y$ обозн. $Y^X$.

\medskip
\noindent
\textbf{Утв.} Если $X, Y$ — конечные, то $|Y^X| = |Y|^{|X|}$

\medskip
\noindent
\textbf{Док-во:} Урновая схема: для каждого из $|X|$ элементов выбираем один из $Y$.

\medskip
\noindent
\textbf{Опр.} Функция $f$ называется инъекцией, если
\[ \forall x,y \in X \quad f(x) = f(y) \Rightarrow x = y \]

\noindent
\textbf{Опр.} Функция $f$ называется сюръекцией, если
\[ \forall y \in Y \quad \exists x \in X : f(x) = y \]

\noindent
\textbf{Опр.} Функция $f$ называется биекцией, если она инъекция и сюръекция.

\bigskip
\noindent
Пусть $|X| = n$, $|Y| = k$.

\vspace{0.4cm}

\begin{minipage}[t]{0.48\textwidth}
\noindent
Количество отображений (функций) $= k^n$ 

\vspace{0.3cm}
\centering
\begin{tikzpicture}[>=stealth, scale=1.3]
    \draw (0,0) ellipse (0.3 and 0.7);
    \node at (0, 0.9) {$n$};
    \fill (0,0.3) circle (1.5pt) coordinate (x1);
    \fill (0,-0.3) circle (1.5pt) coordinate (x2);
    
    \draw (1.8,0) ellipse (0.35 and 0.7);
    \node at (1.8, 0.9) {$k$};
    \fill (1.8,0.4) circle (1.5pt) coordinate (y1);
    \fill (1.8,0.1) circle (1.5pt) coordinate (y2);
    \fill (1.8,-0.2) circle (1.5pt) coordinate (y3);
    \fill (1.8,-0.5) circle (1.5pt) coordinate (y4);

    \draw[->] (x1) -- (y1);
    \draw[->] (x1) -- (y2);
    \draw[->] (x1) -- (y3);
    \draw[->] (x1) -- (y4);
    \node[font=\tiny, fill=white, inner sep=1pt] at (0.9, 0.5) {$k$ способ};

    \draw[->] (x2) -- (y1);
    \draw[->] (x2) -- (y2);
    \draw[->] (x2) -- (y3);
    \draw[->] (x2) -- (y4);
    \node[font=\tiny, fill=white, inner sep=1pt] at (0.9, -0.5) {$k$ способ};
\end{tikzpicture}
\end{minipage}
\hfill
\begin{minipage}[t]{0.48\textwidth}
\noindent
Количество инъекций $= A_k^n$

\vspace{0.3cm}
\centering
\begin{tikzpicture}[>=stealth, scale=1.3]
    \draw (0,0) ellipse (0.4 and 0.8);
    \node at (0, 1.0) {$n$};
    \fill[blue] (0,0.3) circle (1.5pt) coordinate (xi1);
    \fill (0,-0.2) circle (1.5pt) coordinate (xi2);
    
    \draw (2.0,0) ellipse (0.4 and 0.8);
    \node at (2.0, 1.0) {$k$};
    \fill (2.0,0.5) circle (1.5pt) coordinate (yi1);
    \fill[blue] (2.0,0.1) circle (1.5pt) coordinate (yi2);
    \fill (2.0,-0.4) circle (1.5pt) coordinate (yi3);

    \draw[->] (xi1) -- (yi1);
    \draw[->] (xi1) -- (yi2);
    \draw[->] (xi1) -- (yi3);
    \node[font=\tiny, above] at (1.0, 0.4) {$k$ способ};

    \draw[->] (xi2) -- (yi1);
    \draw[->] (xi2) -- (yi3);
    \node[font=\tiny, below, align=center] at (1.0, -0.4) {$k-1$\\способ};

    \draw[->, gray!60] (xi2) -- (yi2);
    \draw[red, thick] (1.15, 0.1) -- (1.35, -0.1);
    \draw[red, thick] (1.35, 0.1) -- (1.15, -0.1);
\end{tikzpicture}
\end{minipage}

\vspace{0.8cm}
\noindent
Количество биекций $= 
\begin{cases} 
0, & \text{если } n \neq k \\ 
n!, & \text{если } n = k 
\end{cases}$

\medskip
\noindent
$\left( \text{Если } n = k\text{, то любая инъекция --- биекция $n! = A_n^n$} \right)$


\medskip
\noindent
Количество сюръекций обозначим $\hat{S}(n, k)$

\medskip
\noindent
Любое отображение --- сюръекция на свой образ
\[ k^n = \sum_{i=0}^{k} C_k^i \cdot \hat{S}(n, i) \]

\medskip
\noindent Число сюръекций в выбранные $i$ элементов --- $\hat{S}(n, i)$

\medskip
\noindent Число сюръекций в какие-то $i$ элементов --- $C_k^i \cdot \hat{S}(n, i)$

\medskip
\noindent Сумма по всем $i$ дает все отображения $k^n$.

\medskip
\noindent \textbf{Замечание} Утверждается, что $\forall \text{ послед. } (f_i)_{i \in \mathbb{N}}, (g_i)_{i \in \mathbb{N}}$
\[ \text{если } f_k = \sum_{i=0}^{k} C_k^i g_i, \text{ то } g_k = \sum_{i=0}^{k} (-1)^{k-i} C_k^i f_i \]

\medskip
\noindent Значит,
\[ \hat{S}(n, k) = \sum_{i=0}^{k} (-1)^{k-i} C_k^i \cdot i^n \]


\noindent
\textbf{Задача.} Есть $n$ книг, из которых $a_1$ книг первого типа, $a_2$ --- второго, \dots, $a_k$ --- $k$-ого типа. ($a_1 + \dots + a_k = n$)

\medskip
\noindent
Различных расстановок: $\dfrac{n!}{a_1! a_2! \dots a_k!}$

\medskip
\noindent
Т.к. при упорядочивании книг каждого типа, которых $a_i!$, мы получим все $n!$ перестановок.

\medskip
\noindent
Числа $\dfrac{n!}{a_1! \dots a_k!}$ называются \textbf{полиномиальными коэффициентами}
\newpage

\medskip
\noindent
\textbf{Обозн.} $P(n; a_1, a_2, \dots, a_k) = C_n^{a_1, \dots, a_k} = \begin{pmatrix} n \\ a_1, \dots, a_k \end{pmatrix}$
\medskip
\[ C_n^{a_1, a_2} = C_n^{a_1}, \text{ т.к. } a_2 = n - a_1 \]

\noindent
\textbf{Посчитаем по-другому:}

\noindent 
Выберем сначала из $n$ ровно $a_1$ позиций для книг первого типа, из оставшихся $n - a_1$ выберем $a_2$ и так далее.
\[ C_n^{a_1, \dots, a_k} = C_n^{a_1} \cdot C_{n-a_1}^{a_2} \cdot C_{n-a_1-a_2}^{a_3} \cdot \dots \cdot C_{n-a_1-\dots-a_{k-1} (=a_k)}^{a_k} \]
Причем можно выбрать $a_i$ в любом порядке --- ответ не изменится.


\vspace{0.4cm}
\noindent
Каждое отображение соответствует упорядоченному разбиению множества $X$ на $k$ частей.

\nopagebreak
\begin{center}
\renewcommand{\arraystretch}{1.2}
\begin{tabular}{l|l}
$f \colon \{1, 2, 3\} \to \{a, b\}$ & $a \ \ \vline \ \ b$ \\
\hline
$f(1) = a$ или $f(1) = b$ & $\varnothing, \{1, 2, 3\}$ \\
$f(2) = a$ или $f(2) = b$ & $\{1\}, \{2, 3\}$ \\
$f(3) = a$ или $f(3) = b$ & $\{2\}, \{1, 3\}$ \\
& $\{3\}, \{1, 2\}$ \\
\hfill $2^3 \text{ способов}$ & $\{1, 2\}, \{3\}$ \\
& $\{2, 3\}, \{1\}$ \\
& $\{1, 3\}, \{2\}$ \\
& $\{1, 2, 3\}, \varnothing$ \\
\end{tabular}
\end{center}

\noindent Если сгруппировать по мощностям частей разбиения получим тождество
\[ k^n = \sum_{\substack{a_1, \dots, a_k \\ a_1 + \dots + a_k = n \\ a_i \ge 0}} C_n^{a_1, \dots, a_k} \]

\noindent Но тогда, поскольку сюръекции --- это в точности разбиения на непустые множества, легко видеть, что
\[ \hat{S}(n, k) = \sum_{\substack{a_1, \dots, a_k \\ a_1 + \dots + a_k = n \\ a_i > 0}} C_n^{a_1, \dots, a_k} \]

\noindent \textbf{Опр.} Число способов разбить множество из n различимых элементов ровно на k непустых неупорядоченных подмножеств называется числом Стирлинга второго рода $S(n,k)$.

\noindent Заметим, что $\hat{S}(n,k)$ --- это число упорядоченных таких же разбиений.

\nopagebreak
\noindent
Поэтому $S(n,k) = \dfrac{\hat{S}(n,k)}{k!}$

\medskip
\noindent
\textbf{Утв.} Числа Стирлинга второго рода удовлетворяет соотношению:
\[ S(n,k) = S(n-1, k-1) + k \cdot S(n-1, k) \]

\noindent
Считаем, что 
$\begin{aligned}[t]
&S(n,0) = 0, n > 0 \\
&S(n,k) = 0, n < k \\
&S(0,0) = 1
\end{aligned}$

\medskip
\noindent
\textbf{Док-во:}
Зафиксируем один элемент \{$x$\}. Если в разбиении есть часть из одного элемента \{$x$\}, то таких разбиений $S(n-1,k-1)$. Посчитаем число оставшихся разбиений.

\medskip
\noindent Распределим все, кроме $x$ --- это $S(n-1, k)$ способов.

\medskip
\noindent
Закинем $x$ в одну из частей $k$ способами

\medskip
\noindent
$\Rightarrow S(n,k) = S(n-1, k-1) + k \cdot S(n-1, k)$
\noindent

$n$ --- предметов, $k$ --- ящиков

\newpage

\nopagebreak
\begin{table}[h!]
\centering
\renewcommand{\arraystretch}{1.8}
\begin{tabular}{|l|l|c|c|c|}
\hline
\textbf{ящики} & \textbf{предметы} & \textbf{\begin{tabular}[c]{@{}c@{}}В ящике\\ сколько угодно\end{tabular}} & \textbf{\begin{tabular}[c]{@{}c@{}}В ящике\\ $\le 1$\end{tabular}} & \textbf{\begin{tabular}[c]{@{}c@{}}В ящике\\ $\ge 1$ предметов\end{tabular}} \\ \hline
разл.          & разл.             & $k^n$                                                                      & $A_k^n$                                                              & $\hat{S}(n,k)$                                                              \\ \hline
разл.          & неразл.           & $\left(\!\binom{k}{n}\!\right)$                                            & $C_k^n$                                                              & $\left(\!\binom{k}{n-k}\!\right) = C_{n-1}^{k-1} = C_{n-k+k-1}^{n-k}$      \\ \hline
неразл.        & разл.             & $B(n,k)$                                                                   & $\begin{matrix} 0, n > k \\ 1, n \le k \end{matrix}$                 & $S(n,k)$                                                                    \\ \hline
\end{tabular}
\end{table}

\vspace{0.4cm}
\noindent
\textbf{Утв.} $B(n,k) = \sum_{i=1}^{k} S(n,i)$ 

\medskip
\noindent
\textbf{Обозн.} В случае, если $n = k$, то $B(n,k)$ обозначается $B_n$ и называется \textbf{числом Белла}.


\section*{\centerline{\Large\textbf{\S \,\, Рекуррентные отношения}}}
\vspace{0.4cm}

\textbf{Задача:} Дана полоска $2 \times n$. Сколькими способами можно ее разбить на доминошки?

\noindent \textbf{Решение:} 

\begin{center}
\begin{tikzpicture}[scale=0.5, every node/.style={font=\small}]
  \draw[step=1, gray!40, very thin] (0,0) grid (8,2);
  \draw[thick] (0,0) rectangle (8,2);
  \node at (4,1) {$A_n$};
  \node at (7.5,1.5) {$\times$};
  \node[left] at (-0.5,1) {все способы};

  \begin{scope}[yshift=-3.5cm]
    \draw[step=1, gray!40, very thin] (0,0) grid (8,2);
    \draw[thick] (0,0) rectangle (8,2);
    \draw[thick, fill=gray!30] (7,0) rectangle (8,2);   
    \draw[dashed] (7,0) -- (7,2);
    \node at (3.5,1) {$A_{n-1}$};
    \node at (7.5,1.5) {$\times$};
    \node[left] at (-0.5,1) {$B_n$: способы с вертикальной};
  \end{scope}

  \begin{scope}[yshift=-7cm]
    \draw[step=1, gray!40, very thin] (0,0) grid (8,2);
    \draw[thick] (0,0) rectangle (8,2);
    \draw[thick, fill=gray!30] (6,1) rectangle (8,2);  
    \draw[thick, fill=gray!30] (6,0) rectangle (8,1);   
    \draw[dashed] (6,0) -- (6,2);
    \node at (3,1) {$A_{n-2}$};
    \node at (7.5,1.5) {$\times$};
    \node[left] at (-0.5,1) {$C_n$: способы с горизонтальной};
  \end{scope}
\end{tikzpicture}
\end{center}

\noindent Тогда $A_n=B_n+C_n$
Заметим, что $B_n=A_{n-1}$, а $C_n=A_{n-2}$, т.к. нижняя угловая клетка однозначно заполняется горизонтальной доминошкой
$A_n=A_{n-1}+A_{n-2}$

\medskip
\noindent \textbf{Опр.} Последовательность $F_1=F_2=1$ и $F_{n+1} = F_n+F_{n+1}$ называются числами Фибоначчи. Чтобы понять как наша задача с ними связана

\begin{tikzpicture}[x=0.8cm, y=0.8cm, every node/.style={font=\scriptsize}]

  \draw (0,0) rectangle (1,2);
  \node at (0.5,2.4) {$2\times1$};
  \node[font=\normalsize] at (2.2,1) {$-\ 1$};
  \node[font=\normalsize] at (4,1) {$=F_2$};

  \begin{scope}[yshift=-2.8cm]
    \draw (0,0) rectangle (2,2);
    \draw (1,0) -- (1,2);
    \node at (1,2.4) {$2\times2$};
    \draw (3.5,0) rectangle (5.5,2);
    \draw (3.5,1) -- (5.5,1);
    \node at (4.5,2.4) {$2\times2$};
    \node[font=\normalsize] at (2.75,1) {,};
    \node[font=\normalsize] at (7,1) {$-\ 2$};
    \node[font=\normalsize] at (8.8,1) {$=F_3$};
  \end{scope}

  \begin{scope}[yshift=-5.6cm]
    \draw (0,0) rectangle (3,2);
    \draw (2,0) -- (2,2);
    \node[font=\normalsize] at (1,1) {?};
    \node at (1.5,2.4) {$2\times3$};
    \node[font=\normalsize] at (1,-0.6) {2};
    \draw (4.5,0) rectangle (7.5,2);
    \draw (5.5,0) -- (5.5,2);
    \draw (5.5,1) -- (7.5,1);
    \node at (6,2.4) {$2\times3$};
    \node[font=\normalsize] at (5.5,-0.6) {1};
    \node[font=\normalsize] at (3.75,1) {,};
    \node[font=\normalsize] at (9,1) {$-\ 3$};
    \node[font=\normalsize] at (10.8,1) {$=F_4$};
  \end{scope}

\end{tikzpicture}

$\Rightarrow A_n = F_{n+1}$

\medskip
\noindent Такой способ задания отношения называется рекуррентным.

\medskip
\noindent\textbf{Опр.} Линейными рекуррентными отношениями называются отношения, в которых предыдущие значения используются лишь в первой степени. 

\noindent \textbf{Пример:}
\begin{align*}
    F_{n+1} = F_n + F_{n-1} & \quad - \text{лин.} \\
    A_n = A_{n-1}^2 & \quad - \text{не лин.} \\
    A_n = A_{n-1} A_{n-2} & \quad - \text{не лин.} \\
    A_n = 3A_{n-1} + 2^n & \quad - \text{лин.}
\end{align*}


\noindent \textbf{Опр.} Рекуррента называется однородной, если все слагаемые одной степени и в ней отсутствует свободный член (слагаемое, не зависящее от членов последовательности).

\noindent \textbf{Пример:}
\begin{align*}
    F_{n+1} = F_n + F_{n-1} & \quad - \text{однород.} \\
    A_n = 3A_{n-1} + 4A_{n-2} + 2^n & \quad - \text{не однород.}
\end{align*}

\noindent \subsection*{Линейные однородные рекурренты}

\noindent \textbf{1)} $a_n = C \cdot a_{n-1}$\\
$a_n = C \cdot a_{n-1} = C \cdot (C \cdot a_{n-2}) = C^2 \cdot a_{n-2} = \dots = C^n \cdot a_0$

\medskip
\noindent \textbf{2)} $a_n = 4a_{n-1} - 3a_{n-2}, \quad a_0 = 2, \ a_1 = 4$

\vspace{0.5em}
\noindent Будем искать решения вида $a_n = \lambda^n$:
\begin{align*}
    \lambda^n &= 4 \cdot \lambda^{n-1} - 3 \cdot \lambda^{n-2} \quad \Big| : \lambda^{n-2} \\
    \lambda^2 &= 4\lambda - 3 \\
    \lambda_{1,2} &= 1; 3 \quad \text{(решим квадратное уравнение)}
\end{align*}

\noindent Если проигнорируем начальные условия, то решениями будут $1^n$ и $3^n$.

\vspace{0.5em}
\noindent Заметим, что решениями также будут $C_1 \cdot 1^n$ и $C_2 \cdot 3^n$.

\noindent Более того, решением будет любая сумма $C_1 \cdot 1^n + C_2 \cdot 3^n$.

\noindent Такое называется \textbf{общим решением}.

\bigskip
\noindent Чтобы найти частное решение, нужно подставить начальные условия:
$$
\begin{cases}
    a_0 = 2 = C_1 \cdot 1^0 + C_2 \cdot 3^0 \\
    a_1 = 4 = C_1 \cdot 1^1 + C_2 \cdot 3^1
\end{cases}
$$
$$ \Rightarrow C_1 = C_2 = 1 \implies a_n = 3^n + 1 $$

\bigskip
\noindent \textbf{3)} $a_n = \alpha_1 \cdot a_{n-1} + \alpha_2 \cdot a_{n-2} + \dots + \alpha_k \cdot a_{n-k}$ \\
$\alpha_1, \dots, \alpha_k$ — коэффициенты

\vspace{0.5em}
\noindent Напишем уравнение:
$$ \lambda^n = \alpha_1 \cdot \lambda^{n-1} + \alpha_2 \cdot \lambda^{n-2} + \dots + \alpha_k \cdot \lambda^{n-k} $$
Поделим на $\lambda^{n-k}$:
$$ \lambda^k = \alpha_1 \cdot \lambda^{k-1} + \dots + \alpha_{k-1} \cdot \lambda + \alpha_k $$
Оно называется \textbf{характеристическим уравнением}.

\medskip
\noindent \textbf{Утв.} Уравнение $n$-ой степени всегда имеет $n$ корней с учётом кратности и комплексных корней. \\
$\left(\text{кратность корня } \lambda \text{ — это максимальное } k \text{ такое, что } P(x) \ \vdots \ (x - \lambda)^k\right)$

\vspace{1.5em}
\noindent Если у нас есть $n$ \underline{различных} корней $\lambda_1, \dots, \lambda_n$, то \textbf{общим решением} называется
$$ C_1 \lambda_1^n + C_2 \lambda_2^n + \dots + C_k \lambda_k^n $$
для констант $C_1, \dots, C_k$

\vspace{1em}
\noindent Поймём, что каждое такое является решением.

\noindent Рассмотрим $a_n = \lambda_i^n$
$$ \lambda_i^n \stackrel{?}{=} \alpha_i \lambda_i^{n-1} + \dots + \alpha_k \lambda_i^{n-k}, \quad \text{что верно так как } \lambda \text{ — его корень} $$

\vspace{0.5em}
\noindent Пусть есть решение $a_n$. Тогда $C \cdot a_n$ — тоже решение.
\begin{align*}
    (C \cdot a_n) &\stackrel{?}{=} \alpha_1 \cdot (C \cdot a_{n-1}) + \alpha_2 \cdot (C \cdot a_{n-2}) + \dots + \alpha_k \cdot (C \cdot a_{n-k}) \quad \Big| : C \\
    a_n &= \alpha_1 a_{n-1} + \dots + \alpha_k a_{n-k} \quad \checkmark
\end{align*}

\noindent Пусть есть 2 решения $a_n$ и $b_n$. Тогда $(a_n + b_n)$ — тоже решение.
$$ (a_n + b_n) \stackrel{?}{=} \alpha_1 (a_{n-1} + b_{n-1}) + \dots + \alpha_k (a_{n-k} + b_{n-k}) $$
$$ \text{Сгруппируем по } a \text{ и по } b \quad \checkmark $$

\vspace{1em}
\noindent Чтобы задать последовательность однозначно, нужно $k$ первых членов. Утверждается, что система с $k$ переменными $(C_1 \dots C_k)$ и $k$ уравнениями (для параметров первых $k$ членов) имеет ровно 1 решение.

\noindent Поэтому мы нашли в точности все решения.

\vspace{1em}
\noindent \textbf{Пример:} $F_n = F_{n-1} + F_{n-2}, \quad F_1 = 1, \ F_2 = 1$
\begin{align*}
    \lambda^2 &= \lambda + 1 \quad \text{— характеристическое уравнение}\\
    \lambda_{1,2} &= \frac{1 \pm \sqrt{5}}{2} \quad \text{— его корни}
\end{align*}

\noindent Тогда общее решение:
$$ F_n = C_1 \cdot \left(\frac{1+\sqrt{5}}{2}\right)^n + C_2 \cdot \left(\frac{1-\sqrt{5}}{2}\right)^n $$

\vspace{0.5em}
\noindent Найдём константы. Часто полезно подставлять 0-ой и 1-ый члены, если они не заданы, то посчитать их:
$$ F_2 = F_1 + F_0 \implies F_0 = 0 $$
$$
\begin{cases}
    0 = F_0 = C_1 \cdot 1 + C_2 \cdot 1 \\
    1 = F_1 = C_1 \cdot \left(\frac{1+\sqrt{5}}{2}\right) + C_2 \cdot \left(\frac{1-\sqrt{5}}{2}\right)
\end{cases}
$$
$$ \Rightarrow C_1 = -C_2 \implies C_1 \cdot \left(\frac{1+\sqrt{5}}{2} - \frac{1-\sqrt{5}}{2}\right) = 1 $$
$$ \Rightarrow C_1 \cdot \sqrt{5} = 1 \implies C_1 = \frac{1}{\sqrt{5}}, \ C_2 = -\frac{1}{\sqrt{5}} $$

$$ F_n = \frac{1}{\sqrt{5}} \left( \left(\frac{1+\sqrt{5}}{2}\right)^n - \left(\frac{1-\sqrt{5}}{2}\right)^n \right) $$
\noindent \textbf{Замеч.} Это называется \textbf{формулой Бине}.

\vspace{1.5em}
\noindent \textbf{4)} $a_n = 10a_{n-1} - 25a_{n-2}, \quad a_0 = 3, \ a_1 = 5$ %(TODO: на лекции $a_1$ было равно 4, но во второй системе, оно уже равно 5, наверно стоит поменять значение $a_1 = 5$)
\begin{align*}
    \lambda^2 &= 10\lambda - 25 \quad \text{— характеристическое уравнение}\\
    (\lambda - 5)^2 &= 0 \implies \lambda_1 = 5 \quad 
\end{align*}
$$
\begin{cases}
    3 = C_1 \cdot 5^0 \\
    5 = C_1 \cdot 5^1
\end{cases} \quad \times
$$
\noindent Рассмотрим рекурренту, у которой корни характеристического уравнения равны $x$ и $y$.

\noindent Тогда $x^n, y^n$ — это общие решения.

\vspace{0.5em}
\noindent Заметим, что $\frac{x^n - y^n}{x - y}$ тоже решение.

\vspace{0.5em}
\noindent Если $x$ и $y$ сближать, то $\frac{x^n - y^n}{x - y}$ превратится в $n \cdot x^{n-1}$.

\vspace{0.5em}
\noindent В нашей задаче будет подходить решение $n \cdot 5^{n-1}$.

\vspace{1em}
\noindent \textbf{Общее уравнение:} $C_1 \cdot 5^n + C_2 \cdot n \cdot 5^{n-1} = C_1 \cdot 5^n + C_2' \cdot n \cdot 5^n$

\vspace{1em}
$$
\begin{cases}
    3 = C_1 \cdot 5^0 + C_2 \cdot 0 \cdot 5^0 \\
    5 = C_1 \cdot 5^1 + C_2 \cdot 1 \cdot 5^1 
\end{cases} \Rightarrow C_1 = 3, \ C_2 = -2
$$

\vspace{1em}
$$ a_n = 3 \cdot 5^n - 2 \cdot n \cdot 5^n $$
\noindent \textbf{5)} $a_n = \alpha_1 a_{n-1} + \dots + \alpha_k a_{n-k}$

\vspace{0.5em}
\noindent У характеристического уравнения корни $\lambda_1, \dots, \lambda_l$ с кратностями $m_1, \dots, m_l$.

\vspace{0.5em}
\noindent Тогда общее решение выглядит:
\begin{align*}
    a_n = \ &C_{1,1} \lambda_1^n + C_{1,2} \cdot n \cdot \lambda_1^n + \dots + C_{1,m_1} \cdot n^{m_1-1} \cdot \lambda_1^n + \dots + C_{l,m_l} \cdot n^{m_l-1} \cdot \lambda_l^n
\end{align*}

\subsection*{Линейные неоднородные уравнения}
\[
a_n = \alpha_1 \cdot a_{n-1} + \dots + \alpha_k \cdot a_{n-k} + f(n)
\]

\[
\text{Общее решение} = \underbrace{\text{Однородное решение}}_{a_n = \alpha_1 a_{n-1} + \dots + \alpha_k a_{n-k}} \quad + \underbrace{\text{Частное решение}}_{b_n = \alpha_1 b_{n-1} + \dots + \alpha_k b_{n-k} + f(n)}
\]

\noindent Если найти какое-то частное решение $b_n$, то $\forall$ однородного решения $a_n$, $b_n + a_n$ тоже будут решением неоднородной рекурренты.

\[
\begin{aligned}
b_n &= \alpha_1 b_{n-1} + \dots + \alpha_k b_{n-k} + f(n) \\
a_n &= \alpha_1 a_{n-1} + \dots + \alpha_k a_{n-k} \\
a_n + b_n &= \alpha_1 (a_{n-1} + b_{n-1}) + \dots + \alpha_k (a_{n-k} + b_{n-k}) + f(n)
\end{aligned}
\]

\bigskip
%лекция от 20.09
\section*{\centerline{\Large\textbf{\S \,\, Дискретная теория вероятностей}}}
\vspace{0.4cm}

$\Omega$ — множество — пространство элементарных исходов.

\noindent \textbf{Пример:} Подбрасывание монетки $\Omega = \{\text{''решка''}, \text{''орел''}\}$

\noindent Бросок кубика $\Omega = \{1, \dots, 6\}$

\noindent $Pr$ — функция $\Omega \rightarrow \mathbb{R}_{+, 0}$, \quad $\sum_{\omega \in \Omega} Pr(\omega) = 1$

\vspace{0.5em}
\noindent Событием называется подмножество $\Omega$, говорят, что событие произошло при элементарном исходе $\omega$, если $\omega$ принадлежит этому подмножеству.

\vspace{0.5em}
\noindent $\mathcal{A}_0$ — это множество всех событий. Из $\mathcal{A}_0$ можно получить $\mathcal{A}$ — множество всех возможных пересечений и объединений событий из $\mathcal{A}_0$.

\vspace{0.5em}
\noindent
$\mathcal{A}$ — алгебра событий.

\medskip
\noindent \textbf{Опр.} Вероятностным пространством называется тройка $(\Omega, \mathcal{A}, Pr)$, где $\Omega, \mathcal{A}, Pr$ описаны выше.

\vspace{0.5em}
\noindent
Мы будем изучать случаи конечного и счётного $\Omega$. В этом случае $\mathcal{A}$ можно считать равным $2^{\Omega}$ (множеством всех подмножеств).

\vspace{0.5em}
\noindent
\textbf{Пример:} Подбрасывание монетки, пока не выпадает орёл.

\noindent
$\Omega = \{1, 2, 3, \dots\}$ — сколько раз подбросили монетку.

\noindent
$Pr(\text{1 подбрасывание}) = \frac{1}{2}$

\noindent
$Pr(\text{2 подбрасывания}) = \frac{1}{4}$

\noindent
$\dots$

\noindent
$\sum_{\omega \in \Omega} Pr(\omega) = 1$

\noindent
В этом случае $\mathcal{A}$ счетно, хотя вероятность всех событий больше 0.

\vspace{0.5em}
\noindent Можно говорить про $Pr(A)$ для события $A$:
\[
Pr(A) = \sum_{\omega \in A} Pr(\omega)
\]

\noindent
\textbf{Свойства:} 
\begin{align*}
& Pr(\varnothing) = 0, \quad Pr(\Omega) = 1 \\
& Pr(A \cup B) = Pr(A) + Pr(B) - Pr(A \cap B)
\end{align*}

\vspace{0.5em}
\noindent
\textbf{Опр.} События $A$ и $B$ называются \textit{несовместными}, если $Pr(A \cap B) = 0$.

$$Pr(A \cup B) = Pr(A) + Pr(B), \text{если $A$ и $B$ несовместны}.$$
$$Pr(\bar{A}) = 1 - Pr(A)$$

\vspace{0.5em}
\noindent
\textbf{Опр.} События $A_1, \dots, A_n$ называются полной группой несовместных событий, если $A_1, \dots, A_n$ --- это разбиение $\Omega$.

\noindent
В этом случае $Pr(A_1) + \dots + Pr(A_n) = 1$.
\\\\
Часто элементарные исходы равновероятны, то есть $Pr(\omega_i) = Pr(\omega_j) \quad \forall i, j$.

\[
Pr(\omega) = \frac{1}{|\Omega|}
\]
\[
Pr(A) = \frac{|A|}{|\Omega|}
\]

\vspace{0.5em}
\noindent
\textbf{Пример:} Бросок кубика
\begin{align*}
& A \text{ — выпадание чётного числа} \\
& B \text{ — выпадание числа } > 3
\end{align*}

\[
Pr(A) = \frac{|A|}{|\Omega|} = \frac{3}{6} = \frac{1}{2}, \quad Pr(B) = \frac{|B|}{|\Omega|} = \frac{3}{6} = \frac{1}{2}
\]

\[
Pr(A \cap B) = \frac{|A \cap B|}{|\Omega|} = \frac{2}{6} = \frac{1}{3}
\]

\noindent
\textbf{Опр.} \textit{Условная вероятность} — это вероятность события $A$ при условии, что событие $B$ уже произошло.

\noindent
\textbf{Обозн.} $Pr(A \mid B)$

\vspace{0.3em}
\noindent
При событии $B$ пространство элементарных исходов сужается с $\Omega$ на $B$.

\noindent
Поэтому:
\[
Pr(A \mid B) = \frac{Pr(A \cap B)}{Pr(B)}
\]
\noindent
\begin{center}
\begin{tikzpicture}[scale=0.7]
    % ===== Левый рисунок: Omega, A, B =====
    \def\ax{-0.6} \def\bx{0.6} \def\r{1.2}

    \draw (0,0) circle (2.6);
    \node at (0,2.1) {$\Omega$};

    % заливка пересечения A \cap B
    \begin{scope}
        \clip (\ax,0) circle (\r);
        \fill[gray!35] (\bx,0) circle (\r);
    \end{scope}
    \draw (\ax,0) circle (\r);
    \draw (\bx,0) circle (\r);

    \node at (-1.3,0.8) {$A$};
    \node at (1.3,0.8) {$B$};
    \node[font=\footnotesize] at (0,0) {$A\cap B$};

    % ===== Стрелка =====
    \draw[-{Stealth[length=2.5mm]}, thick] (3,0) -- (4.2,0);

    % ===== Правый рисунок: только B, внутри дуга границы A =====
    \def\cx{6.3}
    \begin{scope}
        \clip (\cx,0) circle (1.4);
        \fill[gray!35] (\cx-1.2,0) circle (1.2);
        \draw (\cx-1.2,0) circle (1.2); % виден только кусок внутри B
    \end{scope}
    \draw (\cx,0) circle (1.4);
    \node at (\cx+0.7,1.0) {$B$};
    \node[font=\footnotesize] at (\cx-0.5,0) {$A\cap B$};
\end{tikzpicture}
\end{center}

\vspace{1.5em}

\noindent
\textbf{Свойства условной вероятности:}
\begin{enumerate}
    \item $Pr(A \mid A) = 1, \quad Pr(\varnothing \mid A) = 0$
    \item $Pr(A \mid B) = 1$, \quad если $B \subset A$
    \item Если $B_1 \cap B_2 = \varnothing$, то:
    \[
    Pr(A \cap (B_1 \cup B_2)) = Pr((A \cap B_1) \cup (A \cap B_2)) = Pr(A \cap B_1) + Pr(A \cap B_2)
    \]
    и
    \[
    Pr((B_1 \cup B_2) \mid A) = Pr(B_1 \mid A) + Pr(B_2 \mid A)
    \]
\end{enumerate}

\[
Pr(B \mid A) + Pr(\bar{B} \mid A) = 1
\]

\vspace{1em}
\noindent
\textbf{Пример }(предыдущий):
\[
Pr(A \mid B) = \frac{Pr(A \cap B)}{Pr(B)} = \frac{1/3}{1/2} = \frac{2}{3}
\]
\[
Pr(A \mid B) = \frac{|A \cap B|}{|B|} = \frac{2}{3}
\]


% Вывод формулы полной вероятности
\noindent
\begin{align*}
Pr(B) \cdot Pr(A \mid B) &= Pr(A \cap B) \\
&\quad\qquad + \\
Pr(\bar{B}) \cdot Pr(A \mid \bar{B}) &= Pr(A \cap \bar{B}) \\
&\quad\qquad = \\
&\quad\quad Pr(A)
\end{align*}

\noindent
\textbf{Утв. (Формула полной вероятности):} \\
Пусть $A_1, \dots, A_n$ — полная группа несовместных событий. Тогда:
\[
Pr(B) = \sum_{i=1}^{n} Pr(B \mid A_i) \cdot Pr(A_i)
\]

\vspace{0.5em}
\noindent
\textbf{Док-во:} \\
События $A_1 \cap B, A_2 \cap B, \dots, A_n \cap B$ — несовместны.
\[
\Rightarrow \sum_{i=1}^{n} Pr(A_i \cap B) = Pr\left(\bigcup_{i=1}^{n} (A_i \cap B)\right)
\]

\vspace{1.5em}
\noindent
\textbf{Пример.} Завод произвёл 100 деталей. 
60 произвёл первый станок, 25 — второй, 15 — третий. 
Первый делает бракованную деталь с вероятностью $1\%$, второй — $3\%$, третий — $2\%$. 
Какова вероятность брака?

\vspace{0.5em}
\noindent
\textbf{Решение:} \\
Пусть событие $B$ — брак.
\begin{align*}
& A_1 \text{ — деталь 1 мастера} \\
& A_2 \text{ — деталь 2 мастера} \\
& A_3 \text{ — деталь 3 мастера}
\end{align*}

\[
Pr(B) = 0{,}01 \cdot 0{,}6 + 0{,}03 \cdot 0{,}25 + 0{,}02 \cdot 0{,}15 = 0{,}006 + 0{,}0075 + 0{,}003 = 0{,}0165\]

\noindent
\textbf{Понять, у кого больше брака (Формула Байеса):}
\[
Pr(A_1 \mid B) = \frac{0{,}01 \cdot 0{,}6}{Pr(B)}
\]
\[
Pr(A_2 \mid B) = \frac{0{,}03 \cdot 0{,}25}{Pr(B)}
\]
\[
Pr(A_3 \mid B) = \frac{0{,}02 \cdot 0{,}15}{Pr(B)}
\]

\noindent
\textbf{Утв.} 

1.\[
Pr(B \mid A) = \frac{Pr(A \mid B) \cdot Pr(B)}{Pr(A)}
\]

2. Пусть $A_1, \dots, A_n$ — полная группа несовместных событий. Тогда:
\[
Pr(A_i \mid B) = \frac{Pr(A_i) \cdot Pr(B \mid A_i)}{Pr(A_1) \cdot Pr(B \mid A_1) + \dots + Pr(A_n) \cdot Pr(B \mid A_n)} - \textbf{Формула Байеса}
\]

\vspace{0.5em}
\noindent
\textbf{Док-во:} \\
1) $Pr(B \mid A) \cdot Pr(A) = Pr(A \cap B) = Pr(A \mid B) \cdot Pr(B)$ \\
2) В знаменателе формула полной вероятности события $B$.

\medskip
\noindent \textbf{Опр.} События $A$ и $B$ называются \textit{независимыми}, если $Pr(A \mid B) = Pr(A)$.

\vspace{0.5em}
\noindent \textbf{Утв.} События $A$ и $B$ независимы тогда и только тогда, когда $Pr(A \cap B) = Pr(A) \cdot Pr(B)$.

\vspace{0.5em}
\noindent \textbf{Док-во:} $Pr(A \cap B) = Pr(A \mid B) \cdot Pr(B)$

\vspace{0.5em}
\noindent \textbf{Пример} Бросок кубика:
\[
Pr(A \cap B) = \frac{1}{3} \neq \frac{1}{4} = Pr(A) \cdot Pr(B)
\]

\vspace{0.5em}
\noindent Пусть событие $C$ — выпало кратное 3 число.
\[
Pr(A \cap C) = \frac{1}{6} = Pr(A) \cdot Pr(C)
\]
\[
Pr(B \cap C) = \frac{1}{6} = Pr(B) \cdot Pr(C)
\]
\noindent \textbf{Опр.} События $A_1, \dots, A_n$ называют \textit{независимыми в совокупности}, если:
\[
Pr(A_1 \cap \dots \cap A_n) = Pr(A_1) \cdot \dots \cdot Pr(A_n)
\]

\vspace{0.5em}
\noindent \textbf{Пример} Бросаем 2 кубика: \\
\begin{align*}
& A \text{ — выпадение чётного числа на первом} \\
& B \text{ — выпадение чётного на втором} \\
& C \text{ — сумма выпавших чисел чётная}
\end{align*}

\vspace{0.5em}
\noindent
Легко заметить, что:
\begin{align*}
& A, B \text{ — независимы} \\
& B, C \text{ — независимы} \\
& A, C \text{ — независимы}
\end{align*}
Но $A, B, C$ в совокупности зависимы.

{\color{gray} V: $\checkmark$}
\end{document}
